%======================================================================
% L3c_Caratheodory.pdf
%======================================================================
\section{Pfaffovy formy, Carathéodory a $T$ jako integrační faktor}
\strana{L3c-1}
\providecommand{\dbar}{\text{đ}}
\providecommand{\xcancel}[1]{\text{\tikz[baseline=(X.base)]{\node[inner sep=0pt] (X) {$\displaystyle #1$}; \draw[NavyBlue] (X.south west) -- (X.north east);}}}

\begin{definice}[Pfaffova forma]
\kw{Pfaffova forma} -- lin.\ forma v diferenciálech nezávislých proměnných (1-forma)
\[
  \dbar\omega(X_1,\dots,X_k) = \sum_{i=1}^k A_i(X_1,\dots,X_k)\,\dd X_i \qquad (1)
\]
\end{definice}
\begin{itemize}
  \item $\exists\,\omega = \omega(X_1,\dots,X_k)$ taková, že (1) je její úplný diferenciál? \\
  \quad $\Updownarrow$ \\
  \quad \uv{$\dd\omega$ je přímo integrabilní}
  \item pokud ano, potom
  \begin{enumerate}[label=\alph*)]
    \item $A_i = (\nabla\omega)_i \ \Leftrightarrow\ \vec A = \nabla\omega$
    \item $\displaystyle\Big(\pdv{A_i}{X_j}\Big) = \Big(\pdv{A_j}{X_i}\Big)$ \quad \dots\ $\binom k2 = \frac12k(k-1)$ \kw{podmínek integrability} \quad $(i,j = 1,\dots,k)$
  \end{enumerate}
  \item srov.\ 3D: \quad $\vec A = \nabla\omega \ \Rightarrow\ \nabla\times\vec A = \nabla\times(\nabla\omega)\equiv 0$
  \item pokud $\dbar\omega$ není úplný diferenciál, může existovat \kw{integrační faktor} -- funkce $\lambda = \lambda(X_1,\dots,X_k)$ taková, že
  \[ \vzorec{\dd\sigma = \lambda\,\dbar\omega} \]
  je úplný diferenciál funkce $\sigma = \sigma(X_1,\dots,X_k)$ \quad ($\dbar\omega$ integrabilní -- \underline{holonomní})
\end{itemize}

\underline{Pozn:} význam $\sigma$: \quad $\dd\sigma = 0 \Leftrightarrow \dbar\omega = 0$
\begin{itemize}
  \item pokud je tedy $\dbar\omega = \dbar Q$ a $\sigma = S$, potom $S = $ konst.\ definuje \kw{adiabaty} (iso-entropy), po kterých platí $\dbar Q = 0$
\end{itemize}

\begin{itemize}
  \item \underline{podmínky pro existenci $\lambda$} \pozn{v originále napsáno $\mu$}
  \[
    \Big(\pdv{(\lambda A_i)}{X_j}\Big) = \Big(\pdv{(\lambda A_j)}{X_i}\Big)
    \ \Rightarrow\
    \boxed{A_i\pdv{\log\lambda}{X_j} - A_j\pdv{\log\lambda}{X_i} = \pdv{A_j}{X_i} - \pdv{A_i}{X_j} \equiv F_{ij}\neq 0}
  \]
  \begin{itemize}
    \item[--] celkem $\frac12k(k-1)$ parc.\ diferenciálních rovnic pro $\lambda(X_1,\dots,X_k)$
    \item[--] $\lambda$, pokud $\exists$, není určeno jednoznačně: \\
    $\lambda$ je int.\ faktor a $\dd\sigma = \lambda\,\dbar\omega$ $\Rightarrow$ $\bar\lambda = \lambda f(\sigma)$ je také int.\ faktor:
    \[ \bar\lambda\,\dbar\omega = f(\sigma)\,\dd\sigma = \dd\Big(\int^\sigma f(\sigma')\,\dd\sigma'\Big) \]
  \end{itemize}
\end{itemize}
%----------------------------------------------------------------------
\strana{L3c-2}
\subsection*{Holonomní vs.\ anholonomní formy}
\begin{itemize}
  \item \kw{holonomní} $\equiv$ integrovatelné (přímo nebo prostřednictvím $\lambda$)
\end{itemize}
\begin{priklad}[Př.]
\begin{enumerate}[label=\arabic*)]
  \item holonomní \\
  $\dd\omega = x_1\dd x_1 + x_2\dd x_2 + x_3\dd x_3 = 0 \ \Leftrightarrow\ x_1^2 + x_2^2 + x_3^2 = r^2$ \quad (sféra)
  \begin{itemize}[nosep]
    \item pokud se nacházíme na sféře o poloměru $r$, není možné ji opustit po trajektorii splňující $\dd\omega = 0$
  \end{itemize}
  \item anholonomní \\
  $\dbar\omega = x_2\dd x_1 + \dd x_2 + \dd x_3 = 0$ \quad isoplochu nedefinuje
  \begin{enumerate}[label=\alph*)]
    \item ověřte, že forma není integrabilní: $\vec A\cdot\operatorname{rot}\vec A\neq 0$
    \item rozmyslete, že libovolné dva body v $\mathbb{R}^3$ lze spojit po trajektorii $\dbar\omega = 0$ \\
    {[hint: spec.\ je možné se do libovolného bodu dostat z počátku $(0,0,0)$]}
  \end{enumerate}
\end{enumerate}
\end{priklad}

\begin{veta}[$\Rightarrow$ Carathéodoryho teorém]
Vektorové pole $\vec A$ asociované s Pfaffovou formou $\dbar\omega = \vec A\cdot\dd\vec r$ je integrovatelné právě tehdy, pokud v libovolném okolí bodu $P$ příslušného prostoru existují body nedosažitelné po (adiabatické) trajektorii $\dbar\omega = 0$.
\end{veta}
\begin{flushright}\small
DK: $\bullet$ Viz např.\ Luscombe $\leftarrow$ Born, \emph{Natural Philosophy of Cause and Chance}, Oxford, 1949 \\
(str.\ 145, dostupné na archive.org) \\
$\bullet$ Boyling, Commun.\ math.\ Phys.\ \textbf{10}, 52--68 (1968) (dostupné z domény MFF)
\end{flushright}

\begin{itemize}
  \item 2.\ TDZ implikuje integrabilitu $\dbar Q$ s integračním faktorem majícím vlastnosti teploty $\Rightarrow$ lze ekvivalentně zformulovat
\end{itemize}
\begin{veta}[2.\ TDZ podle Carathéodoryho]
V okolí libovolného bodu stavového prostoru existují body adiabaticky nedosažitelné.
\end{veta}
%----------------------------------------------------------------------
\strana{L3c-3}
\subsection*{Integrabilita tepla ve 3D stavovém prostoru}
\begin{itemize}
  \item uvažme dva 1-komp.\ plyny v tep.\ rovnováze
\end{itemize}
\begin{enumerate}[label=\alph*)]
  \item $U(p_A,V_A;p_B,V_B) = U_A(p_A,V_A) + U_B(p_B,V_B)$
  \item \underline{tep.\ rovnováha} $\Rightarrow$ problém má o 1 st.\ volnosti méně \\
  \quad $\vartheta = \phi(p_A,V_A) = \phi(p_B,V_B)$ \quad je empir.\ teplota \\
  \quad $\Rightarrow$ lze volit parametrizaci \fbox{$\vartheta, V_A, V_B$}
\end{enumerate}
\begin{align*}
  \dbar Q &= \dd U_A + p_A\dd V_A + \dd U_B + p_B\dd V_B \\
  &= \Big[\Big(\pdv{U_A}{\vartheta}\Big)_{V_A} + \Big(\pdv{U_B}{\vartheta}\Big)_{V_B}\Big]\dd\vartheta
   + \Big[\Big(\pdv{U_A}{V_A}\Big)_\vartheta + p_A\Big]\dd V_A
   + \Big[\Big(\pdv{U_B}{V_B}\Big)_\vartheta + p_B\Big]\dd V_B
\end{align*}
\begin{enumerate}[label=\alph*),start=3]
  \item \underline{podmínky integrability} \\
  $\dbar\omega = \sum A_i\dd X_i$ je integrabilní $\Leftrightarrow$ $\vec A\cdot\operatorname{rot}\vec A = 0$
  \begin{align*}
    F_i &= (\operatorname{rot}\vec A)_i = \varepsilon_{ijk}\pdv{A_k}{x_j} \\
    F_1 &= \pdv{A_3}{x_2} - \pdv{A_2}{x_3} \equiv 0 \\
    F_2 &= \pdv{A_3}{x_1} - \pdv{A_1}{x_3} = \Big(\frac{\partial^2U_B}{\partial\vartheta\,\partial V_B}\Big) + \Big(\pdv{p_B}{\vartheta}\Big)_{V_B} - \Big(\frac{\partial^2U_B}{\partial V_B\,\partial\vartheta}\Big) = \Big(\pdv{p_B}{\vartheta}\Big)_{V_B} \\
    F_3 &= \pdv{A_1}{x_2} - \pdv{A_2}{x_1} = \Big(\frac{\partial^2U_A}{\partial V_A\,\partial\vartheta}\Big) - \Big(\frac{\partial^2U_A}{\partial\vartheta\,\partial V_A}\Big) - \Big(\pdv{p_A}{\vartheta}\Big)_{V_A} = -\Big(\pdv{p_A}{\vartheta}\Big)_{V_A}
  \end{align*}
\end{enumerate}
\[
  \Rightarrow\ \vec A\cdot\operatorname{rot}\vec A = \Big[\Big(\pdv{U_A}{V_A}\Big)_\vartheta + p_A\Big]\Big(\pdv{p_B}{\vartheta}\Big)_{V_B}
  - \Big[\Big(\pdv{U_B}{V_B}\Big)_\vartheta + p_B\Big]\Big(\pdv{p_A}{\vartheta}\Big)_{V_A}
\]
\begin{itemize}
  \item podmínka integrability dává
  \[
    \vzorec{\frac{\big(\pdv{p_A}{\vartheta}\big)_{V_A}}{\big(\pdv{U_A}{V_A}\big)_\vartheta + p_A}
    = \frac{\big(\pdv{p_B}{\vartheta}\big)_{V_B}}{\big(\pdv{U_B}{V_B}\big)_\vartheta + p_B} = f(\vartheta)}
  \]
  je \underline{univerzální} funkce empirické teploty, neboť nemůže záviset na $p_A,V_A;p_B,V_B$ ani na fyz.\ vlastnostech (stav.\ rovnicích) systémů $A$ nebo $B$ \\
  $\Rightarrow$ \pozor{integrabilita $\dbar Q$ implikuje existenci \underline{absolutní teploty}}
  \item uvažme ideální plyn a jeho stavové rovnice vyjádřené pro obecnou empirickou teplotu:
  \textcolor{NavyBlue}{(je-li $f(\vartheta)$ univerzální, lze k jejímu určení využít libovolný systém)}
  \[
    p = \frac NV g(\vartheta) \quad\&\quad U = cNg(\vartheta)
    \qquad \textcolor{NavyBlue}{\leftarrow\ \text{\small\begin{tabular}{l}to, co experiment dává, je skutečnost,\\ že $g(\vartheta)$ je stejná fce empirické\\ teploty v kalor.\ i termické rovnici\end{tabular}}}
  \]
\end{itemize}
%----------------------------------------------------------------------
\strana{L3c-4}
\[
  \Rightarrow\ \frac{g'(\vartheta)}{g(\vartheta)} = f(\vartheta) \ \Leftrightarrow\ \boxed{f(\vartheta) = \dv{\vartheta}\log g(\vartheta)}
  \quad\dots\ \text{\small vztah absolutní teploty a stav.\ rce id.\ plynu}
\]
\begin{itemize}
  \item spec.\ pokud budeme empirickou teplotu měřit právě s pomocí id.\ plynu, potom
  \[ \vartheta = T \ \&\ g(\vartheta) = T \ \Rightarrow\ \vzorec{f(\vartheta) = \frac1T} \]
  $\to$ přitom účinnost Carnotova cyklu s id.\ plynem je $\eta_C(T_H,T_C) = 1 - \dfrac{T_C}{T_H}$ \\
  \quad $\Rightarrow$ ona univerzální funkce je převrácená absolutní teplota
\end{itemize}
\begin{enumerate}[label=\alph*),start=4]
  \item \underline{integrační faktor k $\dbar Q$}
  \begin{itemize}
    \item je-li $\dbar Q$ integrovatelná forma, co je potom integrační faktor?
  \end{itemize}
  \[
    \dd\omega = \lambda\,\dbar Q = \lambda\sum A_i\dd X_i
    \ \Rightarrow\ A_i\pdv{\log\lambda}{X_j} - A_j\pdv{\log\lambda}{X_i} = \pdv{A_j}{X_i} - \pdv{A_i}{X_j}
  \]
  \begin{itemize}
    \item spec.\ $X_{i=1} = \vartheta$ \& $X_{j=2} = V_A$
    \[
      \Big(\pdv{U_A}{\vartheta} + \pdv{U_B}{\vartheta}\Big)\pdv{\log\lambda}{V_A} - \Big(\pdv{U_A}{V_A} + p_A\Big)\pdv{\log\lambda}{\vartheta}
      = \xcancel{\Big(\frac{\partial^2U_A}{\partial\vartheta\,\partial V_A}\Big)} + \pdv{p_A}{\vartheta} - \xcancel{\frac{\partial^2U_A}{\partial V_A\,\partial\vartheta}}
    \]
    \item předp.\ $\lambda = \lambda(\vartheta)$
    \[
      \Rightarrow\ \pdv{\log\lambda}{\vartheta} = -\frac{\pdv{p_A}{\vartheta}}{\big(\pdv{U_A}{V_A}\big) + p_A} = -f(\vartheta)
      \overset{\text{id.\ plyn}}{=} -\pdv{\log g(\vartheta)}{\vartheta}
    \]
    \[
      \Rightarrow\ \pdv{\log\lambda^{-1}(\vartheta)}{\vartheta} = \pdv{\log g(\vartheta)}{\vartheta}
      \ \Rightarrow\ \lambda(\vartheta)^{-1} = g(\vartheta)
    \]
    $\Rightarrow$ \underline{integračním faktorem je skutečně} \underline{inverzní absolutní teplota}
    \[ \vzorec{\lambda = \frac1T} \]
  \end{itemize}
\end{enumerate}
$\Rightarrow$ Carathéodoryho postulát také vede k definici entropie
\[ \vzorec{\dd S = \frac1T\,\dbar Q} \]
%----------------------------------------------------------------------
\strana{L3c-5}
\begin{priklad}[Př.: plyn magnetických částic]
\[ \dd U = \dbar Q - p\dd V + H\dd M \]
\begin{align*}
  \Rightarrow\ \dbar Q &= \dd U + p\dd V - H\dd M \quad\&\quad U = U(T,V,M) \\
  &= \Big(\pdv{U}{T}\Big)\dd T + \Big[\Big(\pdv{U}{V}\Big) + p\Big]\dd V + \Big[\Big(\pdv{U}{M}\Big) - H\Big]\dd M
\end{align*}
$\Rightarrow$ stejným postupem dostáváme z $\vec A\cdot\operatorname{rot}\vec A = 0$
\[
  \frac{-\big(\pdv{H}{T}\big)_{V,M}}{\big(\pdv{U}{M}\big)_{T,V} - H}
  = \frac{\big(\pdv{p}{T}\big)_{V,M}}{\big(\pdv{U}{V}\big)_{T,M} + p} \overset{(\dagger)}{=} \frac1T
\]
($\dagger$) zde potřebujeme analog rovnice integrability pro $\dd S = \frac1T\dd U - \frac HT\dd M$ \\
\quad a navíc Maxwellovu relaci vycházející z integrability volné energie $\dd F(T,V,M) = -S\dd T - p\dd V + H\dd M$, což zatím neumíme \dots
\end{priklad}
%----------------------------------------------------------------------
\strana{L3c-6}
\subsection*{Integrabilita tepla ve 2D prostoru}
-- triviální tvrzení, ale dává samozřejmě také správný faktor $\lambda$

\textbf{Př:} \quad $\dbar Q = \dd U + p\dd V$ \ \& \ $U = U(p,V)$
\[ \Rightarrow\ \dbar Q = \Big[\Big(\pdv{U}{V}\Big)_p + p\Big]\dd V + \Big(\pdv{U}{p}\Big)_V\dd p \]
\begin{enumerate}[label=\alph*)]
  \item je možné nahlédnout, že $\dbar Q = 0$ jednoznačně definuje adiabatickou trajektorii $p = p(V)$:
  \[
    \dbar Q = 0 \ \Rightarrow\ \dv{p}{V}\Big)_{\text{ad}} = \Big[\Big(\pdv{U}{V}\Big)_p + p\Big]\Big(\pdv{U}{p}\Big)_V^{-1}
  \]
  $\Rightarrow$ obyčejná diferenciální rovnice, řešení existuje za velmi obecných podmínek
  \item podmínka integrability:
  \[
    \Big(\frac{\partial^2U}{\partial p\,\partial V}\Big) + \Big(\pdv{p}{p}\Big)_V \neq \Big(\frac{\partial^2U}{\partial V\,\partial p}\Big)
    \ \Leftrightarrow\ 1\neq 0
  \]
  \item integrační faktor
  \[ \dd\sigma = \lambda\Big[\Big(\pdv{U}{V}\Big)_p + p\Big]\dd V + \lambda\Big(\pdv{U}{p}\Big)_V\dd p \]
  \[
    \Rightarrow\ \Big(\pdv{\lambda}{p}\Big)_V\Big[\Big(\pdv{U}{V}\Big)_p + p\Big] + \xcancel{\lambda\Big(\frac{\partial^2U}{\partial p\,\partial V}\Big)} + \lambda
    = \Big(\pdv{\lambda}{V}\Big)_p\Big(\pdv{U}{p}\Big)_V + \xcancel{\lambda\Big(\frac{\partial^2U}{\partial V\,\partial p}\Big)}
  \]
  \begin{itemize}
    \item vydělím $\big(\pdv{\lambda}{p}\big)_V$: \quad
    $\Big[\Big(\pdv{\lambda}{p}\Big)_V\Big]^{-1} = \Big(\pdv{p}{\lambda}\Big)_V$
    \quad \textcolor{NavyBlue}{$\leftarrow$ viz derivování implicitních funkcí}
  \end{itemize}
  \[
    \Rightarrow\ p + \lambda\Big(\pdv{p}{\lambda}\Big)_V = \Big(\pdv{U}{p}\Big)_V\underbrace{\Big(\pdv{\lambda}{V}\Big)_p\Big(\pdv{p}{\lambda}\Big)_V}_{} - \Big(\pdv{U}{V}\Big)_p
  \]
  \textcolor{NavyBlue}{\small
  $\Big(\pdv{\lambda}{V}\Big)_p\Big(\pdv{p}{\lambda}\Big)_V = \dfrac{\big(\pdv{\lambda}{V}\big)_p}{\big(\pdv{\lambda}{p}\big)_V} = -\Big(\pdv{p}{V}\Big)_\lambda$
  \qquad
  $\pdv{V}\,U\big(V,p(V,\lambda)\big) = \Big(\pdv{U}{V}\Big)_\lambda$}
  \[
    \Rightarrow\ \lambda\Big(\pdv{p}{\lambda}\Big)_V + p = -\Big[\Big(\pdv{U}{V}\Big)_p + \Big(\pdv{U}{p}\Big)_V\Big(\pdv{p}{V}\Big)_\lambda\Big] = -\Big(\pdv{U}{V}\Big)_\lambda
  \]
  \textcolor{Red}{\small ! ještě to není ona známá podmínka integrability $\downarrow$}
  \[
    \Rightarrow\ \Big(\pdv{U}{V}\Big)_\lambda + p = -\lambda\Big(\pdv{p}{\lambda}\Big)_V
    = -\lambda\Big(\pdv{p}{(\lambda^{-1})}\Big)_V\Big(\pdv{(\lambda^{-1})}{\lambda}\Big)_V
    = \lambda^{-1}\Big(\pdv{p}{(\lambda^{-1})}\Big)_V
  \]
  \[
    \Rightarrow\ \vzorec{\Big(\pdv{U}{V}\Big)_{\frac1\lambda} + p = \lambda^{-1}\Big(\pdv{p}{(\lambda^{-1})}\Big)_V
    \ \Leftrightarrow\ \Big(\pdv{U}{V}\Big)_T + p = T\Big(\pdv{p}{T}\Big)_V}
    \quad \textcolor{Red}{\Leftrightarrow\ \lambda = \frac1T}
  \]
\end{enumerate}

